\begin{summarybox}
\small
\par\medskip
\begin{minipage}{\linewidth}
\noindent
\textbf{\textcolor{CSummary}{30 秒速记}}\par\nobreak\smallskip
\[
\boxed{y^2=2px,\ M(x_0,y_0),\ y_0\ne0\quad\Longrightarrow\quad
 k=\frac{p}{y_0}}.
\]\par
\end{minipage}
\par

\par\noindent\textbf{不用死背的推导线：}两点平方作差 $\to$ $(y_1-y_2)(y_1+y_2)=2p(x_1-x_2)$ $\to$ 用 $y_1+y_2=2y_0$ $\to$ $k=p/y_0$。
\par\medskip
\begin{minipage}{\linewidth}
\noindent
\textbf{\textcolor{CSummary}{一张表记住“加减的几何分工”}}\par\nobreak\smallskip
设端点为 $M\pm(u,t)$，$t\ne0$：\par
\end{minipage}
\par

\[
\begin{array}{c|c|c}
\text{曲线}&\text{相减：方向}&\text{相加：长度/存在}\\\hline
\text{圆 }x^2+y^2=R^2&x_0u+y_0t=0&u^2+t^2=R^2-OM^2\\
\text{抛物线 }y^2=2px&y_0t=pu&t^2=2px_0-y_0^2
\end{array}
\]
\par\noindent\textbf{圆：}相减得 $OM\perp AB$，相加直接得到 $OM^2+(AB/2)^2=R^2$ 的\textbf{勾股定理}（圆心处的直径方向任意）。
\par\noindent\textbf{抛物线：}相减得方向 $k=p/y_0$，相加给纵向半跨度及实弦条件；相加\textbf{本身不是勾股}。曲线上同纵标点 $T=(y_0^2/(2p),y_0)$ 与 $M$ 的水平距离满足 $t^2=2|p|\,MT$。
\par\noindent\textbf{真正的半弦长勾股：}
\[
\boxed{AM^2=u^2+t^2
 =(2px_0-y_0^2)\left(1+\frac{y_0^2}{p^2}\right)}.
\]
\par\noindent\textbf{记忆核心：}\textbf{差}去偶项、留下方向；\textbf{和}去奇项、留下展开量；最后\textbf{合}起来才得到实际弦长。
\par\medskip
\begin{minipage}{\linewidth}
\noindent
\textbf{\textcolor{CSummary}{反用时，先做两个检查}}\par\nobreak\smallskip
\begin{enumerate}[leftmargin=*,itemsep=0.2em]
 \item \textbf{有没有弦？} $\boxed{2px_0-y_0^2>0}$ 才有两个不同实交点。
 \item \textbf{能否用有限斜率？} $y_0\ne0$ 用 $k=p/y_0$；$y_0=0$ 且存在实弦时，弦为 $x=x_0$，斜率不存在。
\end{enumerate}\par
\end{minipage}
\par

\par\medskip
\begin{minipage}{\linewidth}
\noindent
\textbf{\textcolor{CSummary}{三种方法怎么选？}}\par\nobreak\smallskip
\begin{itemize}[leftmargin=*,itemsep=0.2em]
 \item \textbf{最快：}点差法，约三行，考场最优。
 \item \textbf{最低门槛：}端点纵坐标设为 $y_0\pm t$，只记“中点对称 + 平方差”。
 \item \textbf{常规：}联立 $y=kx+b$，用韦达定理求 $y_1+y_2$。
\end{itemize}\par
\end{minipage}
\par

\par\medskip
\begin{minipage}{\linewidth}
\noindent
\textbf{\textcolor{CSummary}{三个绝不能混淆的情形}}\par\nobreak\smallskip
\[
\begin{array}{c|c}
 2px_0>y_0^2 & \text{有不同实端点}\\
 2px_0=y_0^2 & \text{退化为切点，不是非退化弦}\\
 2px_0<y_0^2 & \text{无实弦}
\end{array}
\]\par
\end{minipage}
\par

\par\noindent\textbf{最小反例：}$y^2=4x$、$M(0,1)$，虽然形式上 $k=2$，但 $2px_0-y_0^2=-1<0$，没有实弦。
\par\medskip
\begin{minipage}{\linewidth}
\noindent
\textbf{\textcolor{CSummary}{高价值迁移}}\par\nobreak\smallskip
给出\textbf{中点}：先判存在性，再求斜率；给出\textbf{斜率}：先得 $y_0=p/k$，再结合动弦经过的点求 $x_0$。\textbf{$x_0$ 不在斜率公式中，不代表可以任意取值。}\par
\end{minipage}
\par

\end{summarybox}
